% Template for ISBI paper; to be used with:
%          spconf.sty  - ICASSP/ICIP LaTeX style file, and
%          IEEEbib.bst - IEEE bibliography style file.
% --------------------------------------------------------------------------
\documentclass{article}
\usepackage{spconf,amsmath,graphicx}

\usepackage{tikz}
\usetikzlibrary{arrows.meta}

\usepackage{amsmath}
\usepackage{amssymb}

% It's fine to compress itemized lists if you used them in the
% manuscript
\usepackage{enumitem}
\setlist{nosep, leftmargin=14pt}

% Example definitions.
% --------------------
\def\x{{\mathbf x}}
\def\L{{\cal L}}

\title{
REPRODUCING MEANFLOW: A REPLICATION STUDY OF ONE-STEP GENERATIVE MODELING
}

\name{Simon Winther Albertsen, Hjalte Bjoernstrup, Said Djafar Said}

\address{
Department of Computer Science, University of Copenhagen\\
MSc Computer Science\\
Group Assignment 2 (Topics 1--2--3--4), 11th of October 2026
}
%
% For example:
% ------------
%\address{School\\
%	Department\\
%	Address}
%
% Two addresses (uncomment and modify for two-address case).
% ----------------------------------------------------------
%\twoauthors
%  {A. Author-one, B. Author-two\sthanks{Some author footnote.}}
%	{School A-B\\
%	Department A-B\\
%	Address A-B}
%  {C. Author-three, D. Author-four\sthanks{The fourth author performed the work
%	while at ...}}
%	{School C-D\\
%	Department C-D\\
%	Address C-D}
%
% More than two addresses
% -----------------------
% \name{Author Name$^{\star \dagger}$ \qquad Author Name$^{\star}$ \qquad Author Name$^{\dagger}$}
%
% \address{$^{\star}$ Affiliation Number One \\
%     $^{\dagger}$}Affiliation Number Two
%
\begin{document}
%\ninept
%
\maketitle
%
\begin{abstract}
..
\end{abstract}

\begin{keywords}
..
\end{keywords}
%


\section{Introduction}


Diffusion and flow-based generative models produce high-quality images by progressively transforming noise into samples from a learned data distribution. This process typically requires repeated neural network evaluations, making sampling computationally expensive. Reducing the number of evaluations while preserving sample quality is therefore an important objective in generative modelling.

MeanFlow addresses this problem by learning the average velocity over a time interval rather than the instantaneous velocity used in conventional Flow Matching \cite{meanflow}. This representation allows the model to approximate the transformation from noise to an image in a single network evaluation. Its training objective follows from an identity relating average and instantaneous velocities, with the required total derivative computed through a Jacobian-vector product. The method supports training from scratch without a pretrained teacher, distillation, or a progressively adjusted training curriculum.

The effectiveness of this approach depends on whether its reported experimental findings can be reproduced under independently executed training and evaluation procedures. Differences in implementation, training configuration, and computational resources may affect the resulting performance. A reproducibility study must therefore distinguish between reproducing a reported numerical result and recovering the behaviour that supports the underlying methodological claims.

In this report, we investigate selected CIFAR-10 experiments using the authors' released PyTorch implementation. We assess generation quality using Fréchet Inception Distance (FID) and interpret the results in relation to the number of neural network evaluations required for sampling. We document the experimental configuration and any departures from the reference setup, and examine whether our findings support the selected claims of the MeanFlow study.



\section{Conclusion}


% In this assignment, you will select one paper from the curated lists provided under Topic 1 and Topic 2.

% The assignment consists of two main parts:

%     Synopsis of the chosen paper

% Provide a clear description of the main theoretical and methodological contributions of the selected paper in your own words.


%     Discussion of a closely related paper

% Select a second paper outside the curated list that is either a follow-up or a precursor to your chosen paper. Tools such as Google Scholar (https://scholar.google.com/Links to an external site.), Connected Papers (https://www.connectedpapers.com/Links to an external site.), or similar can be useful for identifying such works.

% Your discussion should:

%     Motivate why this second paper was chosen.
%     Explain how the two papers are related (e.g., building upon each other, addressing a common research question).
%     Provide a methodological description of the second paper (including mathematical formulations where relevant).
%     Highlight any methodological improvements or novel aspects compared to the first paper.

% Format and Submission Requirements

%     Use the ISBI LaTeX template and follow all formatting guidelines:
%     Author Instructions (ISBI 2027)Links to an external site.
%     Download Template (ZIP)Links to an external site.
%     Links to an external site.Maximum length: 4 pages (excluding references, figures, and tables).
%     Do not adjust the font size, line spacing, or page margins.
%     Submission: PDF file uploaded via Absalon.

% Additional Notes

%     Illustrations should generally be your own work. It is acceptable to include figures or graphs directly from the papers when relevant (e.g., results), but methodological illustrations/sketches should be original.
%     Think of this assignment as a mini theoretical section of a master's thesis: clear, precise, and written in an academic style. Use mathematical notation where appropriate.
%     Please check the details on the exam format, the use of AI, and the evaluation criteria in the course description.




% References should be produced using the bibtex program from suitable
% BiBTeX files (here: strings, refs, manuals). The IEEEbib.bst bibliography
% style file from IEEE produces unsorted bibliography list.
% ------------------------------------------------------------------------- 

\bibliographystyle{IEEEbib}
\bibliography{strings,refs}

\end{document}
